\begin{codebox}
\codeline{\textcolor{cmtgray}{\#\ Ontology\ Development\ 101\ 七步法走查\ -\ 制造本体}}
\codeline{\textcolor{cmtgray}{\#\ The\ Seven\ Steps\ of\ Ontology\ Development\ 101\ (Noy\ \&\ McGuinness,\ 2001)}}
\codeline{\textcolor{cmtgray}{\#}}
\codeline{\textcolor{cmtgray}{\#\ 以"车间设备本体"为例，完整走一遍七步法}}
\codeblank{}
\codeline{\textcolor{cmtgray}{\#\ ==============================}}
\codeline{\textcolor{cmtgray}{\#\ Step\ 1.\ 确定领域与范围}}
\codeline{\textcolor{cmtgray}{\#\ ==============================}}
\codeline{\textcolor{cmtgray}{\#\ 用四个问题界定：}}
\codeline{\textcolor{cmtgray}{\#\ \ \ 本体覆盖什么领域？\ \ \ \ \ \ \ \ ——\ 离散制造车间的设备与加工能力}}
\codeline{\textcolor{cmtgray}{\#\ \ \ 本体用来做什么？\ \ \ \ \ \ \ \ \ \ ——\ 设备调度、能力匹配、冲突检测}}
\codeline{\textcolor{cmtgray}{\#\ \ \ 本体要回答哪些问题？\ \ \ \ \ \ ——\ 能力问题清单（见\ competency-questions.txt）}}
\codeline{\textcolor{cmtgray}{\#\ \ \ 谁使用和维护本体？\ \ \ \ \ \ \ \ ——\ 工艺工程师建模，调度系统消费}}
\codeblank{}
\codeline{\textcolor{cmtgray}{\#\ 范围控制原则：CQ覆盖不到的概念不建模}}
\codeline{\textcolor{cmtgray}{\#\ （例：本体不需要覆盖财务折旧，除非有对应CQ）}}
\codeblank{}
\codeline{\textcolor{cmtgray}{\#\ ==============================}}
\codeline{\textcolor{cmtgray}{\#\ Step\ 2.\ 考虑复用已有本体}}
\codeline{\textcolor{cmtgray}{\#\ ==============================}}
\codeline{\textcolor{cmtgray}{\#\ 检索顺序（资源详见附录A）：}}
\codeline{\textcolor{cmtgray}{\#\ \ \ 顶层本体：BFO（ISO/IEC\ 21838-2）——\ 物理对象/过程/质量的顶层区分}}
\codeline{\textcolor{cmtgray}{\#\ \ \ 领域本体：IOF（工业本体基础）\ \ \ ——\ 工业过程、设备、能力等核心概念}}
\codeline{\textcolor{cmtgray}{\#\ \ \ 计量本体：QUDT\ \ \ \ \ \ \ \ \ \ \ \ \ \ \ \ \ \ ——\ 单位与量纲（功率kW、长度mm）}}
\codeline{\textcolor{cmtgray}{\#\ \ \ 时间本体：W3C\ Time\ Ontology\ \ \ \ \ ——\ 时间区间与先后关系}}
\codeblank{}
\codeline{\textcolor{cmtgray}{\#\ 复用策略：}}
\codeline{\textcolor{cmtgray}{\#\ \ \ 直接导入（owl:imports）\ \ \ \ \ \ \ \ \ ——\ 适合稳定的小型本体}}
\codeline{\textcolor{cmtgray}{\#\ \ \ 概念映射（rdfs:subClassOf\ 对齐）——\ 适合只需对齐顶层的场景}}
\codeblank{}
\codeline{\textcolor{cmtgray}{\#\ ==============================}}
\codeline{\textcolor{cmtgray}{\#\ Step\ 3.\ 枚举重要术语}}
\codeline{\textcolor{cmtgray}{\#\ ==============================}}
\codeline{\textcolor{cmtgray}{\#\ 不分类、不纠结，先穷举（来自访谈记录与工艺文件）：}}
\codeline{\textcolor{cmtgray}{\#\ \ \ 设备、车床、数控车床、铣床、加工中心、机器人、AGV、}}
\codeline{\textcolor{cmtgray}{\#\ \ \ 材料、铝合金、钛合金、45号钢、}}
\codeline{\textcolor{cmtgray}{\#\ \ \ 工序、车削、铣削、钻孔、热处理、检验、}}
\codeline{\textcolor{cmtgray}{\#\ \ \ 工作单元、产线、车间、}}
\codeline{\textcolor{cmtgray}{\#\ \ \ 功率、转速、精度、利用率、状态、}}
\codeline{\textcolor{cmtgray}{\#\ \ \ 加工、位于、需要、先于、操作……}}
\codeblank{}
\codeline{\textcolor{cmtgray}{\#\ ==============================}}
\codeline{\textcolor{cmtgray}{\#\ Step\ 4.\ 定义类与类层次}}
\codeline{\textcolor{cmtgray}{\#\ ==============================}}
\codeline{\textcolor{cmtgray}{\#\ 三种策略：}}
\codeline{\textcolor{cmtgray}{\#\ \ \ 自顶向下：Equipment\ \(\rightarrow\)\ ProcessingEquipment\ \(\rightarrow\)\ CNCMachine\ \(\rightarrow\)\ CNCLathe}}
\codeline{\textcolor{cmtgray}{\#\ \ \ 自底向上：先有\ CNCLathe/CNCMill，再归纳出\ CNCMachine}}
\codeline{\textcolor{cmtgray}{\#\ \ \ 混合法（推荐）：先定中层骨干概念，向上挂顶层、向下细化}}
\codeblank{}
\codeline{\textcolor{cmtgray}{\#\ 本例类层次（节选）：}}
\codeline{\textcolor{cmtgray}{\#\ Equipment\ \(\sqsubseteq\)\ PhysicalObject}}
\codeline{\textcolor{cmtgray}{\#\ \ \ ProcessingEquipment\ \(\sqsubseteq\)\ Equipment}}
\codeline{\textcolor{cmtgray}{\#\ \ \ \ \ CNCMachine\ \(\sqsubseteq\)\ ProcessingEquipment}}
\codeline{\textcolor{cmtgray}{\#\ \ \ \ \ \ \ CNCLathe\ \(\sqsubseteq\)\ CNCMachine}}
\codeline{\textcolor{cmtgray}{\#\ \ \ \ \ \ \ CNCMillingMachine\ \(\sqsubseteq\)\ CNCMachine}}
\codeline{\textcolor{cmtgray}{\#\ \ \ MeasurementEquipment\ \(\sqsubseteq\)\ Equipment}}
\codeline{\textcolor{cmtgray}{\#\ Material\ \(\sqsubseteq\)\ PhysicalObject}}
\codeline{\textcolor{cmtgray}{\#\ \ \ MetalMaterial\ \(\sqsubseteq\)\ Material}}
\codeline{\textcolor{cmtgray}{\#\ \ \ \ \ AluminumAlloy\ \(\sqsubseteq\)\ MetalMaterial}}
\codeblank{}
\codeline{\textcolor{cmtgray}{\#\ 常见错误自查：}}
\codeline{\textcolor{cmtgray}{\#\ \ \ 单一子类（某类只有一个子类）\(\rightarrow\)\ 层次可能多余}}
\codeline{\textcolor{cmtgray}{\#\ \ \ 类名用复数\ \(\rightarrow\)\ 统一用单数（CNCLathe\ 而非\ CNCLathes）}}
\codeline{\textcolor{cmtgray}{\#\ \ \ 把实例当类（"3号车床"是个体，不是类）}}
\codeblank{}
\codeline{\textcolor{cmtgray}{\#\ ==============================}}
\codeline{\textcolor{cmtgray}{\#\ Step\ 5.\ 定义属性}}
\codeline{\textcolor{cmtgray}{\#\ ==============================}}
\codeline{\textcolor{cmtgray}{\#\ 对象属性（连接两个个体）：}}
\codeline{\textcolor{cmtgray}{\#\ \ \ canProcess\ \ \ :\ Equipment\ \(\rightarrow\)\ Material}}
\codeline{\textcolor{cmtgray}{\#\ \ \ locatedIn\ \ \ \ :\ Equipment\ \(\rightarrow\)\ WorkCell}}
\codeline{\textcolor{cmtgray}{\#\ \ \ precedes\ \ \ \ \ :\ Process\ \(\rightarrow\)\ Process（传递属性）}}
\codeline{\textcolor{cmtgray}{\#\ \ \ operates\ \ \ \ \ :\ Worker\ \(\rightarrow\)\ Equipment}}
\codeblank{}
\codeline{\textcolor{cmtgray}{\#\ 数据属性（连接个体与字面量）：}}
\codeline{\textcolor{cmtgray}{\#\ \ \ hasPower\ \ \ \ \ \ \ \ :\ Equipment\ \(\rightarrow\)\ xsd:float（单位kW）}}
\codeline{\textcolor{cmtgray}{\#\ \ \ hasSerialNumber\ :\ Equipment\ \(\rightarrow\)\ xsd:string}}
\codeline{\textcolor{cmtgray}{\#\ \ \ hasStatus\ \ \ \ \ \ \ :\ Equipment\ \(\rightarrow\)\ \{Running,\ Idle,\ Maintenance\}}}
\codeblank{}
\codeline{\textcolor{cmtgray}{\#\ ==============================}}
\codeline{\textcolor{cmtgray}{\#\ Step\ 6.\ 定义属性约束（Facets）}}
\codeline{\textcolor{cmtgray}{\#\ ==============================}}
\codeline{\textcolor{cmtgray}{\#\ 基数约束：每台设备恰有一个序列号}}
\codeline{\textcolor{cmtgray}{\#\ \ \ Equipment\ \(\sqsubseteq\)\ =1\ hasSerialNumber}}
\codeline{\textcolor{cmtgray}{\#\ 值域约束：功率非负}}
\codeline{\textcolor{cmtgray}{\#\ \ \ hasPower\ 的取值\ \(\ge\)\ 0}}
\codeline{\textcolor{cmtgray}{\#\ 全称约束：数控机床的控制器都是数控系统}}
\codeline{\textcolor{cmtgray}{\#\ \ \ CNCMachine\ \(\sqsubseteq\)\ \(\forall\)hasControl.CNCControl}}
\codeline{\textcolor{cmtgray}{\#\ 不相交：Equipment\ \(\sqcap\)\ Material\ \(\sqsubseteq\)\ \(\bot\)}}
\codeblank{}
\codeline{\textcolor{cmtgray}{\#\ ==============================}}
\codeline{\textcolor{cmtgray}{\#\ Step\ 7.\ 创建实例}}
\codeline{\textcolor{cmtgray}{\#\ ==============================}}
\codeline{\textcolor{cmtgray}{\#\ Lathe\_003\ :\ CNCLathe}}
\codeline{\textcolor{cmtgray}{\#\ \ \ hasSerialNumber\ "EQ-2023-0117"}}
\codeline{\textcolor{cmtgray}{\#\ \ \ hasPower\ 15.5}}
\codeline{\textcolor{cmtgray}{\#\ \ \ locatedIn\ WorkCell\_A}}
\codeline{\textcolor{cmtgray}{\#\ \ \ canProcess\ AluminumAlloy\_6061}}
\codeblank{}
\codeline{\textcolor{cmtgray}{\#\ 实例数据通常批量来自台账/MES（导入管道见第7章）}}
\codeblank{}
\codeline{\textcolor{cmtgray}{\#\ ==============================}}
\codeline{\textcolor{cmtgray}{\#\ 迭代提醒}}
\codeline{\textcolor{cmtgray}{\#\ ==============================}}
\codeline{\textcolor{cmtgray}{\#\ 七步不是瀑布而是循环：创建实例时发现缺属性\ \(\rightarrow\)\ 回到Step\ 5；}}
\codeline{\textcolor{cmtgray}{\#\ CQ验收不通过\ \(\rightarrow\)\ 回到Step\ 4\ 调整层次。}}
\codeline{\textcolor{cmtgray}{\#\ 每轮迭代后运行推理器做一致性检查（工具见第4章Protégé）}}
\end{codebox}
